% Part: incompleteness
% Chapter: incompleteness-provability
% Section: introduction

\documentclass[../../../include/open-logic-section]{subfiles}

\begin{document}

\olfileid[fr]{inc}{inp}{int}

\olsection{Introduction}

Hilbert estimait qu'un système d'axiomes décrivant
une structure mathématique, par exemple les entiers
naturels, est insuffisant s'il ne permet pas de
démontrer tous les énoncés vrais sur cette structure.
Cet objectif, joint à son intérêt ultérieur pour les
systèmes formels de déduction, suggère qu'il voulait
garantir que les systèmes employés pour raisonner
sur les entiers naturels soient non seulement
cohérents, mais aussi \emph{complets} : tout énoncé
de leur langage, ou sa négation, doit être
!!{derivable}. Le premier théorème d'incomplétude
de Gödel montre qu'il n'existe aucun système
d'axiomes de l'arithmétique à la fois complet,
cohérent et !!{axiomatizable}. Plus généralement,
aucune théorie mathématique cohérente,
!!{axiomatizable} et « suffisamment forte »
n'est complète.

Un objectif encore plus important de Hilbert,
au cœur de son programme de justification des
mathématiques modernes (« classiques »), était de
trouver des preuves finitaires de cohérence pour les
systèmes formels représentant le raisonnement
classique. À cet égard, le second théorème
d'incomplétude de Gödel porta un coup bien plus
grave à son programme. Comme le premier, on peut
l'énoncer de façon encore vague : en substance,
aucune théorie cohérente et !!{axiomatizable} de l'arithmétique,
suffisamment forte, ne peut démontrer sa propre cohérence. Pour l'argument présenté ici,
« suffisamment forte » désignera une théorie
un peu plus puissante que~$\Th{Q}$.

L'idée de la démonstration originale de Gödel se
retrouve dans le paradoxe d'Épiménide. Crétois
lui-même, Épiménide affirmait que tous les Crétois
sont des menteurs; la forme plus directe du paradoxe
est « cette phrase est fausse ». En remplaçant la
vérité par la !!{derivability}, Gödel a construit
un !!{sentence} qui affirme indirectement qu'il
n'est pas lui-même !!{derivable}. Si cet
!!{sentence} était !!{derivable}, la théorie serait
incohérente. Gödel a également montré que, dans le
système d'axiomes qu'il considérait, la négation de
cet !!{sentence} n'était pas !!{derivable}.
Pour cette seconde partie, il devait supposer que
la théorie~$\Th{T}$ était « $\omega$-cohérente ».
Cette propriété est plus forte que la simple
cohérence.\footnote{Toute théorie $\omega$-cohérente
est cohérente, mais la réciproque est fausse.}
Quelques années après Gödel, Rosser a montré que
la seule cohérence de~$\Th{T}$ suffisait.

La première difficulté consiste à construire un
!!{sentence} qui se rapporte à lui-même. Pour
toute !!{formula}~$!A$ du langage de~$\Th{Q}$,
notons $\gn{!A}$ le numéral correspondant au
nombre~$\Gn{!A}$. Distinguons les niveaux : $!A$
est une !!{formula} du langage de~$\Th{Q}$,
$\Gn{!A}$ est un entier naturel, et $\gn{!A}$
est un \emph{terme} du langage de~$\Th{Q}$.
Chaque !!{formula}~$!A$ possède ainsi un
\emph{nom}, $\gn{!A}$, dans ce langage. Ce codage
permet à une !!{formula} de « parler » d'autres
!!{formula}s. Le résultat suivant est le lemme
du point fixe.

\begin{lem}
Soit $\Th{T}$ une théorie qui étend~$\Th{Q}$,
et soit $!B(x)$ une !!{formula} dont la seule
variable libre est~$x$. Il existe alors un
!!{sentence}~$!A$ tel que
$\Th{T}\Proves !A\liff !B(\gn{!A})$.
\end{lem}

Le lemme dit que, pour toute propriété~$!B(x)$,
on peut construire un !!{sentence}~$!A$ qui
affirme « $!B(x)$ est vraie de moi » et que
$\Th{T}$ « connaît » cette équivalence.

Comment construire un tel !!{sentence}?
Considérons cette version du paradoxe d'Épiménide,
due à Quine :
\begin{quote}
« donne une phrase fausse quand on le fait précéder
de sa propre citation » donne une phrase fausse quand
on le fait précéder de sa propre citation.
\end{quote}
Cette phrase n'est pas directement autoréférentielle.
Elle énonce une propriété du groupe de mots cité,
comme les phrases suivantes :
\begin{enumerate}
\item « Robert » est un joli prénom.
\item « J'ai couru. » est une phrase courte.
\item « A trois mots » a trois mots.
\end{enumerate}
Mais que se passe-t-il si l'on place, devant le
groupe de mots « donne une phrase fausse quand on
le fait précéder de sa propre citation », sa
propre citation? On obtient la phrase initiale.
Elle affirme donc qu'elle est fausse.

\end{document}
